Every entry satisfies $f(f(f(x)))=x$. Direct factorization gives
\[
f^{\circ3}(x)-x=(x^2-5x+2)G(x)H(x),
\]
where
\[
G=x^3-7x^2+10x+7,\qquad H=x^3-4x^2+x+1.
\]
Roots of $x^2-5x+2$ are fixed points and cannot occur in a distinct three-cycle. Both cubics have three distinct real roots: $G$ changes sign on each of $(-1,0),(3,4),(4,5)$, while $H$ does so on $(-1,0),(0,1),(3,4)$.

The identities
\[
G(f(x))=G(x)(x^3-5x^2+2x+1),
\]
\[
H(f(x))=H(x)(x^3-8x^2+17x-5)
\]
show that $f$ preserves the root set of each cubic. Neither cubic shares a root with $x^2-5x+2$: their remainders upon division by it are $11-2x$ and $4x-1$, respectively, and the corresponding potential common roots do not satisfy the quadratic. Consequently each three-root set is a genuine three-cycle.

By Vieta, the products of the roots of $G,H$ are $-7,-1$. Hence the complete answer is $\boxed{K\in\{1,7\}}$.
